\minitoc 

% =======================================================================================================
% Introduction 

Intéressons nous maintenant aux problèmes de programmation linéaire. Même si le nom peut prêter à confusion, 
il s'agit en réalité de problèmes d'optimisation d'une \emph{fonction objectif} en respectant 
certaines \emph{contraintes}. Le caractère linéaire de tels problèmes vient du fait que 
nous ne traiterons que des fonctions objectifs et des contraintes linéaires. 

Ces problèmes de programmation linéaires peuvent être résolus en temps polynomial, notament grâce à l'algorithme 
du simple. Les problème de flots maximaux, vus précédement, peuvent s'exprimer comme des problèmes de programmation 
linéaire. 

% =======================================================================================================
% Programmation Linéaire

\section{Programmation Linéaire}

\subsection{Généralités}

\begin{definition}[Programme linéaire]
    Un \emph{programme linéaire} ou \emph{problème de programmation linéaire} est un problème d'optimisation 
    de $n \in \N$ variables et $ m \in \N$ contraintes de la forme : 
        \begin{mini*}|l|
            {}{z = \sum_{j=1}^{n} c_j x_j }{}{}
            % \addConstraint{g_{j}(x)}{\le 0, j = 1, \hdots, m}
            \addConstraint{ \sum_{j=1}^{n} a_{ij} x_j}{\preceq b_i \quad \forall i \in \llbracket 1, m \rrbracket }
            \addConstraint{        x }{ \succeq 0 \quad \forall j \in \llbracket 1, n \rrbracket }
        \end{mini*}
    où $ \preceq \in \{ \leqslant, \geqslant, =\}$. Un programme linéaire est donc un simple problème d'opmisation 
    sur un polyèdre convexe. Dans le cas définit plus haut, on dit que le programme est sous forme générale. Différentes 
    formes d'un même problème existent. 
\end{definition}

\begin{example}
    Voyons dès à présent un exemple de programme linéaire. 
    \begin{quote}
        \emph{
            Un barman possède 30 litres de jus d'orange, 24 litres de pamplemousse et 18 litres de jus de 
            framboise. Il peut préparer 2 types de cocktails : 
            \begin{itemize}
                \item Le "sunrise", vendu 80 euros le baril de 10 litres composé de 5 litres de de jus d'orange, de 
                2 litres de jus de pamplemouse, d'un litre de jus de framboise et de 2 litres d'eau. 
                \item Le "swing", vendu 60 euros le baril et composé de 3 litres de jus de pamplemouse, 3 litres de jus d'orange, 
                3 litres de jus de framboise et d'un libre d'eau. 
            \end{itemize}
            Il souhaite préparer la plus importante quantité des deux cocktails pour maximiser son profit. 
        }
    \end{quote}
    Pour résoudre ce problème, nous pouvons construire un programme linéaire de deux variables $x$ et $y$ représentant chacunes 
    le nombre de litres d'un des deux cocktails à produire. La fonction objectif, décrivant le prix de 
    revente de la marchandise est donc : 
        $$ z = 80 x + 60 y $$ 
    Et nous avons les contraintes suivantes : 
        \begin{align*}
            5x + 3y \leqslant 30 & \quad \text{(stock de jus d'orange)} \\
            2x + 3y \leqslant 24 & \quad \text{(stock de jus de pamplemouse)} \\
            x + 3y \leqslant 18 & \quad \text{(stock de jus de framboise)} \\ 
            x,z \geqslant 0 & 
        \end{align*}
    On peut donc formuler le problème suivant : 
        \begin{maxi*}|l|
            {}{z = 80 x + 60 y }{}{}
            \addConstraint{ 5x + 3y }{ \leqslant 30 }
            \addConstraint{ 2x + 3y }{ \leqslant 24 }
            \addConstraint{ x + 3y }{ \leqslant 18 }
            \addConstraint{ x,y }{ \geqslant 0 }
        \end{maxi*}
\end{example}

\begin{proposition}
    Comme dit dans la définition d'un programme linéaire, il peut exister différentes formes d'un même problème en 
    fonction de la méthode de résolution utilisée. Le caractère linéaire de ces programmes nous permet de 
    passer facilement d'une forme à une autre. Nous définissons la forme canonique d'un programme linéaire comme suit :

        \begin{figure}[h!]
            \centering
            \begin{maxi*}|l|
                {}{z = \sum_{j=1}^{n} c_j x_j }{}{}
                \addConstraint{ \sum_{j=1}^{n} a_{ij} x_j }{ \preceq b_i, \quad \forall i \in \llbracket 1, m \rrbracket }
                \addConstraint{ x_j }{ \succeq 0, \quad \forall j \in \llbracket 1, n \rrbracket }
            \end{maxi*}
            \caption{Formulation canonique d’un programme linéaire}
        \end{figure}

        \begin{center}
            \begin{minipage}{0.60\textwidth}
                \centering
                \begin{maxi*}|l|
                    {}{z = \sum_{j=1}^{n} c_j x_j }{}{}
                    \addConstraint{ \sum_{j=1}^{n} a_{ij} x_j }{ = b_i, \quad \forall i \in \llbracket 1, m \rrbracket }
                    \addConstraint{ x_j }{ \succeq 0, \quad \forall j \in \llbracket 1, n \rrbracket }
                \end{maxi*}
                \captionof{figure}{Formulation standard d’un programme linéaire de maximisation}
            \end{minipage}
            \hfill 
            \begin{minipage}{0.35\textwidth}
                Nous pouvons ensuite passer à la formulation dite « standard » où les contraintes d’ordre sont remplacées 
                par des contraintes d’égalités. Il faut donc rajouter un certain nombre de variables à notre programme. 
            \end{minipage}
        \end{center}
    
    
    Enfin, nous avons la forme matricielle, qui nous permettra d'utiliser des algorithmes itératifs pour résoudre ce genre de 
    problèmes : 


    \begin{figure}[h!]
        \centering
        \begin{maxi*}|l|
            {}{z = CX }{}{}
            \addConstraint{ AX }{ = B }
            \addConstraint{ X }{ \geqslant 0_{\R^n} }
        \end{maxi*}
        \caption{Formulation matricielle d’un programme linéaire de maximisation}
    \end{figure}

    où $C \in \R^n, A \in \mathcal{M}_{m, n}(\R)$ et $B \in \R^m$.  
\end{proposition}

\begin{example}
    Pour le problème de notre barman, nous pouvons ainsi représenter les différentes formes de notre problème : 

        \begin{center}
            \begin{minipage}{0.45\textwidth}
                \begin{maxi*}|l|
                    {}{z = 80 x + 60 y }{}{}
                    \addConstraint{ 5x + 3y }{ \leqslant 30 }
                    \addConstraint{ 2x + 3y }{ \leqslant 24 }
                    \addConstraint{ x + 3y }{ \leqslant 18 }
                    \addConstraint{ x,y }{ \geqslant 0 }
                \end{maxi*}
                \captionof{figure}{Forme canonique du problème du barman}
            \end{minipage}
            \hfill 
            \begin{minipage}{0.45\textwidth}
                \begin{maxi*}|l|
                    {}{z = 80 x + 60 y }{}{}
                    \addConstraint{ 5x + 3y + e_1 }{ = 30 }
                    \addConstraint{ 2x + 3y + e_2 }{ = 24 }
                    \addConstraint{ x + 3y + e_3 }{ = 18 }
                    \addConstraint{ x,y, e_1, e_2, e_3}{ \geqslant 0 }
                \end{maxi*}
                \captionof{figure}{Forme standard du problème du barman}
            \end{minipage}
        \end{center}
        Et enfin : 
            \begin{maxi*}|l|
                {}{
                    z = \underbrace{(80, 60, 0, 0, 0)}_{C} \times 
                    \overbrace{
                        \begin{pmatrix}
                            x \\ 
                            y \\ 
                            e_1 \\ 
                            e_2 \\ 
                            e_3 
                        \end{pmatrix}
                    }^{X}
                }{}{}
                \addConstraint{ 
                    \underbrace{
                    \begin{pmatrix}
                        5 & 3 & 1 & 0 & 0 \\ 
                        2 & 3 & 0 & 1 & 0 \\ 
                        1 & 3 & 0 & 0 & 1 
                    \end{pmatrix}
                    }_{A}
                    \times 
                    \underbrace{
                        \begin{pmatrix}
                            x \\ 
                            y \\ 
                            e_1 \\ 
                            e_2 \\ 
                            e_3 
                        \end{pmatrix}
                    }_{X}
                }{  =
                    \underbrace{
                        \begin{pmatrix}
                            30 \\ 
                            24 \\ 
                            18 
                        \end{pmatrix}
                    }_{B}
                }
            \end{maxi*}
\end{example}

Maintenant que nous avons défini formellement un programme linéaire, passons à l'étude de ses solutions. 

\begin{definition}[Solutions d'un programme linéaire]
    Soit $P$ un programme linéaire sous forme matricielle standard. Le vecteur $X = (x_1, \dots, x_n)$ est : 
    \begin{itemize}
        \item \emph{une solution de $P$} si $A \; ^tX = B$ 
        \item \emph{une solution possible de $P$} si $A \; ^tX = B$ et $X \geqslant 0_{\R^n}$
        \item \emph{une solution optimale de $P$} si c'est une solution possible et qu'elle 
        maximise $z$ parmi l'ensemble des solutions possibles. 
    \end{itemize}
\end{definition}

\begin{proposition}
    Comme énoncé en introduction, on peut représenter un problème de maximisation de flot dans un réseau par 
    un programme linéaire. Soit $R = (X, E, c)$ un réseau tel que $|E| = n$, avec $s, t \in X$ la source et le puits. 
    On note $c(x,y)$ la capacité de l’arête entre $x$ et $y$ dans le réseau. 
    On cherche une application 
        \[  \varphi : E \longrightarrow \mathbb{R}, \quad (x,y) \longmapsto f(x,y) \]
    qui maximise le flot entre la source $s$ et le puits $t$. 

    Plus formellement, si l’on ordonne les arêtes $E$ sous la forme $E = \{e_1, \dots, e_n\}$ et que l’on note 
    $\varphi(E) = (f_1, \dots, f_n) \in \mathbb{R}^n$, alors on cherche à maximiser 
        \[ z = \sum_{(x,y) \in v^-(s)} \varphi(x,y) \]
    sous les contraintes :
        \[ 0 \leqslant f(x,y) \leqslant c(x,y), \quad \forall (x,y) \in E \]
    ainsi que les contraintes de conservation du flot pour tout sommet $v \in X \setminus \{s, t\}$ :
        \[ \sum_{(u,v) \in v^-(v)} f(u,v) = \sum_{(v,w) \in v^+(v)} f(v,w).\]
\end{proposition}

\subsection{Solution Géométrique}

La recherche d'une solution d'un programme linéaire revient à trouver une droite de la forme $ z = \sum_{i=1}^{n} x_i c_i $
pour laquelle $z$ est maximal tout en respectant les contraintes du programme. 
En petite dimension (2 variables ou moins) il est possible de tracer ce problème et de le résoudre géométriquement. 
Détaillons cela autour d'un exemple. 


\begin{example}[Solution Géométrique]
    Reprenons le problème de notre barman énoncé comme suit :
    \begin{maxi*}|l|
        {}{z = 80 x + 60 y }{}{}
        \addConstraint{ 5x + 3y }{ \leqslant 30 }
        \addConstraint{ 2x + 3y }{ \leqslant 24 }
        \addConstraint{ x + 3y }{ \leqslant 18 }
        \addConstraint{ x,y }{ \geqslant 0 }
    \end{maxi*}

    On peut tracer les courbes des contraintes pour identifier la zone contenant les solutions possibles :

    \centering
    \begin{tikzpicture}
        \begin{axis}[
            axis lines=middle,
            xlabel={$x$},
            ylabel={$y$},
            xmin=0, xmax=8,
            ymin=0, ymax=8,
            grid=major,
            width=10cm,
            height=8cm,
            legend style={at={(0.97,0.97)}, anchor=north east},
            clip=false
        ]

        % === Contraintes ===
        \addplot[name path=C1, thick, domain=0:6, blue] { (30 - 5*x)/3 };
        \addlegendentry{$5x + 3y = 30$}

        \addplot[name path=C2, thick, domain=0:8, red] { (24 - 2*x)/3 };
        \addlegendentry{$2x + 3y = 24$}

        \addplot[name path=C3, thick, domain=0:8, orange] { (18 - x)/3 };
        \addlegendentry{$x + 3y = 18$}
        
        % === Solution ===
        \addplot[name path=C4, dashed, domain=0:8, black] { (540 - 80*x)/60 };
        \addlegendentry{$540 = 80x + 60y$}

        % === Zone faisable ===
        \addplot[fill=green!30, fill opacity=0.4]
            coordinates {(0,0) (0,6) (3,5) (6,0)} -- cycle;

        % === Points d'intersection ===
        \addplot[mark=*] coordinates {(0,6)}; 
        \addplot[mark=*] coordinates {(6,0)}; 
        \addplot[mark=*] coordinates {(3,5)}; 
        \addplot[mark=*] coordinates {(0,0)}; 

        % === Flèche de la fonction objectif ===
        \addplot[->, thick, black!70] coordinates {(0,0) (3,5)};
        \node[black!70, right] at (axis cs:4,3) {$z = 80x + 60y$ ↑};

        \end{axis}
    \end{tikzpicture}

    \captionof{figure}{Espace des solutions réalisables du programme linéaire}

    Le barman peut donc produire 80 litres de "sunrise" et 3 litres de "swing" pour un profit total de 540 euros. 
\end{example}

\begin{remark}
    En pratique, il sera très rare de pouvoir résoudre un tel problème graphiquement car le nombre de variable sera 
    bien supérieur à 2. 
\end{remark}


\subsection{Algorithme du Simplexe}

Comme dit précédement, le caractère linéaire des contraintes d'un programme linéaire nous permettent 
d'en déduire que l'espace des solutions possible est un polyèdre. D'où le théorème suivant. 

\begin{theorem}[Solutions Possibles]
    L'ensemble $S$ des solutions possibels d'un programme linéaire est convexe et est un polyèdre non borné 
    ou un polytope. Si $S$ est non vide et borné, alors il existe au moins un vecteur de $S$ pour lequel 
    le maximum de la fonction objectif est atteint.  
\end{theorem}

\begin{remark}
    Ce théorème nous permet de distinguer plusieurs cas pour l'ensemble des solutions : 
    \begin{itemize}
        \item Le programme linéaire a au moins une solution : 
            \begin{itemize}
                \item soit elle est unique 
                \item soit il en existe une infinité (i.e le maximum est atteint en plusieurs vecteurs) on parle 
                alors de cas dégénéré. 
            \end{itemize}
        \item Le programme linéaire n'a pas de solution, lorsque les contraintes ne sont pas compatibles, 
        l'ensemble des solutions est donc vide. 
        \item Le programme linéaire n'a pas d'optimum (i.e l'ensemble des solutions n'est pas borné).
    \end{itemize}
\end{remark}

\begin{proposition}[Résolution d’un programme linéaire par le simplexe]
    Nous avons précédemment résolu nos programmes linéaires graphiquement. 
    Nous allons voir que nous pouvons aussi les résoudre \emph{algébriquement}, 
    en utilisant l'\emph{algorithme du simplexe}.
    Soit le programme suivant : 

    \begin{maxi*}|l|
        {}{z = x + 2y }{}{}
        \addConstraint{ 4x + 2y }{ \leqslant 300 }
        \addConstraint{ 5x }{ \leqslant 310 }
        \addConstraint{ 2x + 4y }{ \leqslant 200 }
        \addConstraint{ 3y }{ \leqslant 100 }
        \addConstraint{ x, y }{ \geqslant 0 }
    \end{maxi*}

    On passe ce programme sous \emph{forme standard} en introduisant des variables d’écart 
    $e_1, e_2, e_3, e_4 \geqslant 0$ :
    
    \begin{maxi*}|l|
        {}{z = x + 2y}{}{}
        \addConstraint{4x + 2y + e_1}{ = 300} 
        \addConstraint{5x + e_2}{ = 310}
        \addConstraint{2x + 4y + e_3}{ = 200}
        \addConstraint{3y + e_4}{ = 100}
        \addConstraint{x, y, e_1, e_2, e_3, e_4}{ \geqslant 0 }
    \end{maxi*}

    Le vecteur $(x, y, e_1, e_2, e_3, e_4) = (0, 0, 300, 310, 200, 100)$ 
    est une \emph{solution de base réalisable} : toutes les contraintes sont vérifiées et les variables sont positives.
    L’idée du simplexe est de :
    \begin{enumerate}
        \item Choisir une \emph{base initiale} (ici, les variables d’écart $e_i$).
        \item Exprimer les variables hors base ($x, y$) en fonction des variables de base.
        \item \emph{Augmenter la valeur de la fonction objectif $z$} en remplaçant une variable de base par une variable hors base
              qui améliore $z$.
        \item Répéter jusqu’à ce qu’aucune amélioration ne soit possible : la solution est alors optimale.
    \end{enumerate}

    Formellement, le simplexe cherche à écrire le système sous la forme :
    \[
        \begin{cases}
            z = C + \sum_i c_i x_i \\
            x_B = b - A x_N
        \end{cases}
        \quad \text{où } x_B \geqslant 0, \; x_N \geqslant 0
    \]
    À chaque itération :
    \begin{itemize}
        \item une variable hors base entre dans la base (celle qui augmente le plus $z$),
        \item une variable de base quitte la base (celle qui devient nulle en premier).
    \end{itemize}
    Lorsque tous les coefficients $c_i \leqslant 0$, 
    aucune amélioration n’est possible et la solution actuelle est \emph{optimale} : 
    la valeur optimale est alors $z = C$.
\end{proposition}

\begin{example}[Résolution algébrique du programme par le simplexe]
    Considérons le programme linéaire suivant :
    \begin{maxi*}|l|
        {}{z = x + 2y}{}{}
        \addConstraint{4x + 2y}{\leqslant 300}
        \addConstraint{5x}{\leqslant 310}
        \addConstraint{2x + 4y}{\leqslant 200}
        \addConstraint{3y}{\leqslant 100}
        \addConstraint{x, y}{\geqslant 0}
    \end{maxi*}

    \begin{enumerate}
        \item \textbf{Mise sous forme standard : } On introduit une variable d’écart pour chaque contrainte :
                \[
                    \begin{cases}
                        4x + 2y + e_1 = 300 \\
                        5x + e_2 = 310 \\
                        2x + 4y + e_3 = 200 \\
                        3y + e_4 = 100
                    \end{cases}
                    \quad \text{avec } e_i \geqslant 0.
                \]
            La fonction objectif s’écrit :
                \[
                    z = x + 2y.
                \]

            Une solution de base réalisable évidente est obtenue en prenant \(x = y = 0\), ce qui donne :
                \[
                    (e_1, e_2, e_3, e_4) = (300, 310, 200, 100),
                    \quad z = 0.
                \]
            C’est notre {solution de base initiale}.  
        
        \item \textbf{Recherche d’une direction d’amélioration : } Comme \(z = x + 2y\), augmenter \(y\) 
        est plus rentable que \(x\) (le coefficient 2 est supérieur à 1).  
        On va donc chercher à \textbf{augmenter \(y\)} tant que c’est possible sans violer les contraintes.
        On part du système :
            \[
                \begin{cases}
                    e_1 = 300 - 2y - 4x \\
                    e_2 = 310 - 5x \\
                    e_3 = 200 - 2x - 4y \\
                    e_4 = 100 - 3y
                \end{cases}
            \]
        Comme on garde \(x=0\) au départ, ces équations deviennent :
            \[
                \begin{cases}
                    e_1 = 300 - 2y \\
                    e_2 = 310 \\
                    e_3 = 200 - 4y \\
                    e_4 = 100 - 3y
                \end{cases}
            \]
        Pour rester faisable, il faut \(e_i \geq 0\).  
        On obtient donc :
            \[
                y \leq 150, \quad y \leq 50, \quad y \leq 25, \quad y \leq 33{,}3.
            \]
        La contrainte la plus restrictive est \(y \leq 25\) (provenant de \(e_3 = 200 - 4y\)).  
        On choisit donc \(y = 25\) et \(e_3 = 0\).  
        Cela définit une nouvelle base : \(\{e_1, e_2, e_4, y\}\).
        
        \item \textbf{Amélioration suivante : } Pour cette valeur \(y = 25\), les variables d’écart valent :
            \[
                e_1 = 300 - 2(25) = 250, \quad
                e_2 = 310, \quad
                e_3 = 0, \quad
                e_4 = 100 - 3(25) = 25.
            \]
        La fonction objectif vaut alors :
            \[
                z = x + 2y = 2 \times 25 = 50.
            \]
        On peut encore améliorer \(z\) en augmentant \(x\), tant que toutes les contraintes restent valides.
        En remplaçant \(y = 25\), on obtient :
            \[
                \begin{cases}
                    e_1 = 300 - 2(25) - 4x = 250 - 4x \\
                    e_2 = 310 - 5x \\
                    e_3 = 200 - 2x - 4(25) = 100 - 2x \\
                    e_4 = 100 - 3(25) = 25
                \end{cases}
            \]
        Les contraintes donnent :
            \[
                x \leq 62{,}5, \quad x \leq 62, \quad x \leq 50.
            \]
        Le plus restrictif est \(x \leq 50\) (provenant de \(e_3 \geq 0\)).  
        On prend donc \(x = 50\) et \(e_3 = 0\).

        \item \textbf{Calcul de la solution optimale : } On obtient :
            \[
                \begin{cases}
                    x = 50, \quad y = 25, \\
                    e_1 = 300 - 4(50) - 2(25) = 100, \\
                    e_2 = 310 - 5(50) = 60, \\
                    e_3 = 0, \\
                    e_4 = 100 - 3(25) = 25.
                \end{cases}
            \]
        La fonction objectif vaut :
            \[
                z = x + 2y = 50 + 2 \times 25 = 100.
            \]
        Si on essaye d’augmenter \(y\) ou \(x\) davantage, on viole immédiatement une contrainte.  
        La solution est donc \textbf{optimale}.
                
        \item \textbf{Conclusion :} Le programme atteint son maximum en :
            \[
                \boxed{x^* = 50, \quad y^* = 25, \quad z_{\max} = 100.}
            \]
        Cette résolution illustre la logique du simplexe : on part d’une solution réalisable triviale,
        puis on améliore la fonction objectif pas à pas en déplaçant la base vers les sommets adjacents
        du polytope faisable, jusqu’à atteindre un sommet où toute augmentation violerait une contrainte.
    \end{enumerate}
\end{example}


\subsection{Dualité}

\begin{definition}[Programme Primal]
    Le \emph{problème primal} (ou primaire) est la forme standard d’un programme linéaire :
        
        \begin{maxi*}
            {}{z = \; ^tc x}{}{}
            \addConstraint{Ax}{\leqslant b}
            \addConstraint{x}{\geqslant 0}
        \end{maxi*}

    où :
    \begin{itemize}
        \item $x \in \mathbb{R}^n$ est le vecteur des variables de décision ;
        \item $c \in \mathbb{R}^n$ est le vecteur des coefficients de la fonction objectif ;
        \item $A \in \mathbb{R}^{m \times n}$ est la matrice des coefficients des contraintes ;
        \item $b \in \mathbb{R}^m$ est le vecteur des bornes du système d’inégalités.
    \end{itemize}
\end{definition}

\begin{definition}[Problème Dual]
    À tout problème primal est associé un \emph{problème dual} défini par :
        
        \begin{mini*}|l|
            {}{w = \; ^tb y}{}{}
            \addConstraint{^tA y}{ \geqslant c}
            \addConstraint{y}{\geqslant 0}
        \end{mini*}

    où :
    \begin{itemize}
        \item $y \in \mathbb{R}^m$ est le vecteur des variables duales ;
        \item chaque contrainte du primal correspond à une variable du dual ;
        \item chaque variable du primal correspond à une contrainte du dual.
    \end{itemize}
\end{definition}

\begin{remark}[Interprétation]
    Le problème \emph{primal} cherche à \emph{maximiser un profit} ou une valeur objective en respectant des 
    contraintes de ressources. Le problème \emph{dual} cherche à \emph{minimiser le coût} de ces ressources tout 
    en garantissant que la valeur de chaque ressource couvre la contribution de chaque variable du primal.
    Autrement dit :
    \begin{quote}
        \emph{Le primal cherche {ce qu’on peut obtenir au maximum} avec des ressources limitées,  
        tandis que le dual cherche {le coût minimal} que doivent avoir ces ressources pour que le système reste équilibré.}
    \end{quote}
\end{remark}

\begin{example}
    En reprenant l'exemple de notre barman, le programme primal cherche à maximiser le profit du barman 
    en fixant le prix des barrils $c_j$ de cocktails $j$ qu'il produit en utilisant au plus $b_i$ litres 
    de jus $i$. Les variables $x_j$ représentent alors le nombre de barrils de cocktail $j$. 

    Ici, le programme dual cherche à minimiser le profit total obtenu grâce aux stocks $b_i$ de chaque jus 
    $i$ si le barman vent chaque cocktail $j$ pour un prix minimal de $c_j$. 
    Les variables $y_i$ représentent donc le profit que le barman réalise par litre de jus $i$. 

    Voici les expressions mathématiques de ces deux programmes :  

    \begin{center}
        \begin{minipage}{0.45\textwidth}
            \begin{maxi*}|l|
                {}{z = 80 x + 60 y }{}{}
                \addConstraint{ 5x + 3y }{ \leqslant 30 }
                \addConstraint{ 2x + 3y }{ \leqslant 24 }
                \addConstraint{ x + 3y }{ \leqslant 18 }
                \addConstraint{ x,y }{ \geqslant 0 }
            \end{maxi*}
            \captionof{figure}{Programme primal du barman}
        \end{minipage}
        \hfill
        \begin{minipage}{0.45\textwidth}
            \begin{mini*}|l|
                {}{w = 20 y_1 + 24 y_2 + 18 y_3}{}{}
                \addConstraint{5 y_1 + 2 y_2 + y_2}{ \geqslant 80}
                \addConstraint{3y_1 + 3y_2 + 3y_3}{\geqslant 60}
                \addConstraint{y_1, y_2, y_3}{\geqslant 0}
            \end{mini*}
            \captionof{figure}{Programme dual du barman}
        \end{minipage}
    \end{center}

    La solution optimale du primal est $(x_1, x_2) = (3, 5)$ pour un profit maximal de $540$. En revanche 
    la solution du dual est $(y_1, y_2, y_3) = (15, 0, 5)$ pour un profit de $540$ aussi. 
\end{example}

Il existe un lien direct entres les solutions du dual et celles du primal. Elles sont décrites par les théorèmes 
de dualité. 

\begin{theorem}[Dualité Faible]
    La solution optimale du programme primal est inférieure ou égale à celle du dual. 
\end{theorem}

\begin{theorem}[Dualité Forte]
    Nous avons les deux implications suivantes : 
    \begin{enumerate}
        \item Si l'un des deux programmes (primal ou dual) admet une solution optimale finie, alors l'autre aussi. 
        De plus, la valeur optimale pour des deux programmes est la même. 
        \item Si l'ensemble des solutions de l'un des deux programmes est non borné , alors celui de l'autre est vide. 
    \end{enumerate}
\end{theorem}

\begin{remark}
    Le théorème de coupure minimale/flot maximal dans un réseau est l'équivalent du théorème de dualité forte ici. 
    En effet, le problème de coupure minimale dans un réseau est l'exact dual du problème de flot maximal dans le 
    même réseau. 
\end{remark}

\begin{example}[Utilité de la dualité]
    Considérons le programme linéaire suivant : 
        
        \begin{maxi*}|l|
            {}{x + 3x + 6y + z}{}{}
            \addConstraint{w + x + 5y + 6z}{ \leqslant 20}
            \addConstraint{w + 2x + y + 3z}{ \leqslant 30}
            \addConstraint{x,y,w,z}{\geqslant 0}
        \end{maxi*}
    
    Ce programme ne peut pas être résolu graphiquement du fait du trop grand nombre de 
    variables. En passant au dual, nous aurons un programme linéaire à deux variables et 4 contraintes que l'on 
    pourra résoudre graphiquement. Attention, la résolution du dual nous donnera seulement une idée de la 
    valeur de la solution maximale. 
\end{example}

\begin{remark}
    Ainsi, passer au problème dual peut être avantageux dans plusieurs cas : 
    \begin{itemize}
        \item Le dual peut être plus simple à résoudre (moins de variables ou moins de contraintes). 
        \item Il peut exister une solution optimale évidente pour le dual mais pas pour le primal. 
        \item La valeur de n'importe quelle solution du dual nous donne une borne supérieure pour 
        pour la solution du primal. 
        \item La solution optimale du dual nous donne des informations sur la solution optimale du primal. 
    \end{itemize}
\end{remark}

% =======================================================================================================
% Programmation Linéaire en nombres entiers

\section{Programmation Linéaire en Nombres Entiers}

Dans cette section, nous allons traiter un nouveau type de problèmes : les \emph{programes linéaires en nombres entiers}. 
Vous l'aurez compris, il s'agit maintenant de considérer des programmes linéaires à variables entières. Nous allons donc 
une méthode pour les résoudre ainsi que quelques techniques de modélisation. Commençons tout d'abord 
par un exemple introductif. 

\begin{example}
    Un commerçant souhaite remplir son sac à dos d'au plus $n$ objets de poids et valeurs respectifs $ p_i \in \N$ 
    et $ v_i \in \N$. Or le sac du commerçant ne peut contenir qu'un poids maximal $C \in \N$. 
    Le commerçant veut choisir les objets à prendre de manière à maximiser la valeur emportée. 
    On modélisera donc le problème de la façon suivante : 
        
        \begin{maxi*}|l|
            {}{V = \sum_{i=1}^{n} v_i \times x_i }{}{}
            \addConstraint{ \sum_{i=0}^{n} p_i x_i }{\leqslant C} 
            \addConstraint{\forall i \in \llbracket 1, n \rrbracket, \; x_i \in \N}{}
        \end{maxi*}
    Ici le vecteur $(x_1, \dots, x_n) \in \{0,1\}^n$ représentera les objets que l'on prend. En effet, $x_i = 0$ 
    si le commerçant ne prend pas l'objet $i$ et $1$ sinon. 

    Nous avons donc affaire à un programme linéaire. Or sa résolution pourrait non conduire à une solution comportant 
    des valeurs $x_i \not \in \{0,1\}^n$ alors que l'on considère que l'on ne peut pas prendre une partie d'un objet 
    (ex : $3/4$ de l'object $i$). Nous devons donc chercher une nouvelle méthode pour résoudre ce genre de programmes. 
\end{example}

\subsection{Définition et relaxation}

\begin{definition}[Programme linéaire en nombres entiers]
    Un \emph{programme linéaire en nombres entiers} est un programme linéaire comme définit précédement où un certain nombre 
    de variables sont à valeurs entières. Plus formellement, la forme strandard d'un tel programme est de la forme : 
        
        \begin{maxi*}|l|
            {}{\sum_{j=1}^{n} c_j x_j}{}{}
            \addConstraint{\sum_{j=1}^{n} a_{ij} x_j}{ = b_i \quad \forall i \in \llbracket 1, m \rrbracket}
            \addConstraint{x_j}{\geqslant 0 \quad \forall j \in \llbracket 1, n \rrbracket}
            \addConstraint{x_j}{\in \N \quad \forall j \in \llbracket 1, n \rrbracket}
        \end{maxi*}
    où les constantes $c_j, a_{ij}$ et $b_i$ sont entières. 
\end{definition}

\begin{definition}[Relaxation linéaire]
    Soit un programme linéaire en nombres entiers (PLNE) :

        \begin{maxi*}|l|
            {}{\sum_{j=1}^{n} c_j x_j }{}{}
            \addConstraint{\sum_{j=1}^{n} a_{ij} x_j}{= b_i \quad \forall i \in \llbracket 1, m \rrbracket}
            \addConstraint{x_j \in \N \quad \forall j \in \llbracket 1, n \rrbracket}
        \end{maxi*}

    La \emph{relaxation linéaire} de ce programme consiste à supprimer les contraintes de la forme : 
        $$ x_j \in \N \quad \forall j \in \llbracket 1,n \rrbracket $$ 
    On obtient ainsi un programme linéaire classique, plus facile à résoudre. 
\end{definition}

\begin{example}
    Considérons le programme linéaire suivant que nous allons résoudre graphiquement :

        \begin{maxi*}|l|
            {}{z = 2x + 3y}{}{}
            \addConstraint{5 x + 8y}{\leqslant 40}
            \addConstraint{2x-y}{\leqslant 8}
            \addConstraint{x,y}{\geqslant 0}
        \end{maxi*}

        \centering
        \begin{tikzpicture}
            \begin{axis}[
                axis lines=middle, 
                xlabel={$x$}, 
                ylabel={$y$}, 
                xmin= 0, xmax=10,
                ymin=0, ymax=8, 
                grid=major, 
                height=8cm, 
                width=10cm, 
                legend style={at={(0.97,0.97)}, anchor=north east},
                clip=false
            ]
                % === CONTRAINTES === 
                \addplot[name path=C1, thick, domain=-1:8, blue] { (40 - 5*x)/8 };
                \addlegendentry{$5x + 8y \leqslant 40$}

                \addplot[name path=C2, thick, domain=3.8:8, red] {2*x - 8}; 
                \addlegendentry{$2x - y \leqslant 8$}

                % === SOLUTION === 
                \addplot[name path=S, dashed, domain=2:6, gray] {11.81 - 2 * x}; 
                \addlegendentry{$2x + y = 11.81$}


                % === ZONE FAISABLE === 
                \addplot[fill=green!30, fill opacity=0.4]
                    coordinates {(0,0) (4,0) (4.95,1.9) (0,5)} --cycle; 
            \end{axis}
        \end{tikzpicture}
    
    Graphiquement, nous trouvons une solution optimale de valeur $11.81$. 

\end{example}

\subsection{Algorithme "Branch and Bound"}

\begin{theorem}[Lien entre PLNE et relaxation]
    Soit un programme linéaire en nombres entiers (PLNE) et sa relaxation linéaire correspondante.  
    Notons $z^*_\text{int}$ la valeur optimale du PLNE et $z^*_\text{relax}$ celle de la relaxation. Alors :  
    \[
        z^*_\text{relax} \geqslant z^*_\text{int} \quad \text{(en cas de maximisation)}, \qquad
        z^*_\text{relax} \leqslant z^*_\text{int} \quad \text{(en cas de minimisation)}.
    \]
    Autrement dit, la solution optimale de la relaxation fournit une \emph{borne supérieure} (ou \emph{borne inférieure}) pour le PLNE initial.
\end{theorem}

La méthode \emph{Branch and Bound} est une technique de résolution des programmes linéaires en nombres entiers 
(PLNE) qui repose sur l'utilisation des relaxations linéaires pour obtenir des bornes sur la solution optimale.
Cet algorithme construit un arbre de recherche binaire dont la racine est le 
problème initial $\Pi$. L'algorithme s'arrête lorsque tous les nœuds feuilles de l'arbre sont \emph{fathomés} 
(c'est-à-dire qu'ils n'ont plus besoin d'être explorés).

\begin{definition}[Nœud fathomé]
    Un nœud est dit \emph{fathomé} si au moins une des conditions suivantes est vérifiée :
    \begin{enumerate}
        \item La relaxation linéaire est \emph{infaisable}, c'est-à-dire qu'aucune solution réelle ne satisfait les contraintes.
        \item La valeur optimale de la relaxation linéaire est inférieure ou égale à la meilleure solution entière connue.
        \item La solution optimale de la relaxation linéaire est entière, et est donc également solution optimale du programme entier.
    \end{enumerate}
\end{definition}

\begin{algorithm}[H]
\caption{Branch and Bound pour un programme linéaire en nombres entiers}
\begin{algorithmic}[1]
\State Résoudre la relaxation linéaire du problème $\Pi$.
\If{le nœud est fathomé}
    \State arrêter l'exploration de ce nœud.
\Else
    \State Identifier une variable $x_i$ fractionnaire dans la solution $x^*$.
    \State Créer deux sous-problèmes en ajoutant respectivement les contraintes :
    \[
        x_i \leqslant \lfloor x^*_i \rfloor \quad \text{et} \quad x_i \geqslant \lceil x^*_i \rceil
    \]
    \State Appeler \textbf{Branch and Bound} récursivement sur ces sous-problèmes.
\EndIf
\end{algorithmic}
\end{algorithm}

\begin{remark}
    L'algorithme explore donc l'arbre de recherche en combinant deux techniques :
    \begin{itemize}
        \item Le \emph{branching}, qui divise le problème en sous-problèmes plus simples.
        \item Le \emph{bounding}, qui utilise la relaxation linéaire pour élaguer les sous-problèmes 
              ne pouvant pas améliorer la meilleure solution entière connue.
    \end{itemize}
\end{remark}
